\problemname{Prime Sums}

\begin{figure}
\centering
\begin{verbatim}
 163
 547
+829
----
1539
\end{verbatim}
\caption{The three three-digit numbers are all primes and the digits 1..9 are each used exactly once in the three numbers}
\end{figure}

In the figure you see an ordinary addition of three three-digit numbers. What is special about the numbers is that they are all \emph{prime numbers}. Furthermore, each of the nine digits 1..9 appears exactly once in the three numbers.

Write a program that determines how many such additions can be formed. The program should also determine the \emph{highest} and \emph{lowest} values of these sums.

Note that all additions containing the same three numbers are considered the same addition.

\section*{Input}
This problem has no input.

\section*{Output}
Output three integers: the number of additions, the lowest value of a sum, and the highest value of a sum.
